\documentclass[12pt,a4paper]{ctexart}

\usepackage[margin=1.55cm,top=1.35cm,bottom=1.45cm]{geometry}
\usepackage{amsmath,amssymb}
\usepackage{tikz}
\usepackage{xcolor}
\usepackage{enumitem}

\definecolor{Ink}{RGB}{25,34,48}
\definecolor{Accent}{RGB}{34,100,170}
\definecolor{SoftLine}{RGB}{213,224,236}
\definecolor{SoftFill}{RGB}{247,250,253}

\pagestyle{empty}

\setlength{\parindent}{0pt}
\setlength{\parskip}{0.45em}
\setlist[itemize]{leftmargin=1.6em,itemsep=0.2em,topsep=0.2em}
\setlist[enumerate]{leftmargin=1.8em,itemsep=0.45em,topsep=0.25em}

\newcommand{\lessonTitle}[2]{
  {\LARGE\bfseries\color{Ink} #1}\hfill
  {\normalsize\color{Accent} #2}
  \vspace{0.55em}
  {\color{SoftLine}\hrule height 0.8pt}
  \vspace{0.8em}
}

\newcommand{\questionBox}[2]{
  \begin{tikzpicture}
    \node[
      draw=SoftLine,
      fill=SoftFill,
      line width=0.7pt,
      inner sep=10pt,
      text width=\dimexpr\linewidth-22pt\relax,
      align=left
    ] {
      {\large\bfseries\color{Accent} #1}\par\vspace{0.45em}
      #2
    };
  \end{tikzpicture}
}

\newcommand{\writeArea}[1]{
  \vspace{0.7em}
  \begin{tikzpicture}
    \draw[SoftLine,line width=0.7pt] (0,0) rectangle (\linewidth,#1);
    \foreach \y in {0.9,1.8,...,#1} {
      \draw[SoftLine,line width=0.35pt] (0.35,\y) -- (\linewidth-0.35cm,\y);
    }
  \end{tikzpicture}
}

\begin{document}

\lessonTitle{清爽竞赛讲义风}{A4 竖版 · 题目与留白}

\questionBox{第1题：概念与基础计算}{
已知函数 $f(x)=x^2-2ax+1$。
\begin{enumerate}
  \item 写出函数图象的对称轴。
  \item 当 $x\in[0,3]$ 时，讨论 $f(x)$ 的最小值与参数 $a$ 的关系。
\end{enumerate}
}

\writeArea{15.2}

\newpage
\lessonTitle{清爽竞赛讲义风}{A4 竖版 · 大题一页}

\questionBox{第2题：函数零点与参数}{
设函数 $g(x)=\ln x-kx+1$，其中 $k$ 为实数。
\begin{enumerate}
  \item 当 $k=1$ 时，判断 $g(x)$ 在 $(0,+\infty)$ 上的零点个数。
  \item 若函数 $g(x)$ 在 $(0,+\infty)$ 上恰有两个零点，求 $k$ 的取值范围。
\end{enumerate}
}

\vspace{0.5em}
\begin{center}
\begin{tikzpicture}[scale=0.95]
  \draw[->,line width=0.7pt,Ink] (-0.2,0) -- (7.2,0) node[right] {$x$};
  \draw[->,line width=0.7pt,Ink] (0,-2.0) -- (0,3.2) node[above] {$y$};
  \draw[Accent,line width=1.0pt,domain=0.25:6.6,samples=120] plot (\x,{ln(\x)+1-0.55*\x});
  \draw[SoftLine,line width=0.7pt] (1,0.08) -- (1,-0.08) node[below=4pt] {$1$};
  \node[Accent] at (4.95,1.2) {$y=\ln x-kx+1$};
\end{tikzpicture}
\end{center}

\writeArea{10.6}

\newpage
\lessonTitle{清爽竞赛讲义风}{A4 竖版 · 小题智能混排}

\questionBox{第3题：同类小题组}{
\begin{enumerate}
  \item 已知 $\vec a=(2,-1)$，$\vec b=(m,3)$，若 $\vec a\perp \vec b$，求 $m$。
  \item 已知复数 $z=1+2i$，求 $z\overline z$。
  \item 正方体 $ABCD-A_1B_1C_1D_1$ 中，判断直线 $AC$ 与平面 $BDD_1B_1$ 的位置关系。
\end{enumerate}
}

\writeArea{15.0}

\end{document}
